% INF-221 Tarea 1 2026-2 | Tomás San Martín | Rol 202473565-9
Equipo: Intel(R) Core(TM) i5-9300H CPU @ 2.40GHz, 24 GB RAM, Windows 11 con WSL2 (Ubuntu); compilador
g++ 13.3.0 con flags \texttt{-O2 -std=c++17} (en todos los algoritmos). 

Se empleó el nivel de optimización \texttt{-O2} por tres razones. En primera instancia, porque nuestro objetivo
es medir y estudiar el comportamiento de los algoritmos en la práctica. En segundo lugar, por factibilidad experimental, ya que, de lo contrario,
en las pruebas de Strassen, el tiempo de ejecución aumentaría considerablemente. Y en tercer lugar, al medir en un
entorno igual para todos (mismas flags), las diferencias de las curvas las marcan los algoritmos y no el compilador.

Los casos de prueba se generaron con los scripts provistos por el
curso: ordenamiento con $N=\{10,10^3,10^5,10^7\}$, $T=\{$ascendente,
descendente, aleatorio$\}$, $D=\{$D1,D7$\}$, $M=\{$a,b,c$\}$; matrices con
$N=\{2^4,2^6,2^8,2^{10}\}$, $T=\{$dispersa, diagonal, densa$\}$,
$D=\{$D0,D10$\}$, $M=\{$a,b,c$\}$ (72 casos por problema). Como limitaciones en
el desarrollo de la tarea tenemos efectos de caché al momento de comparar teórica y experimentalmente los
algoritmos, y la resolución del muestreo de memoria.